% GENERATED FILE. Edit canonical lessons, then run npm run generate:books.
% canonical-source-hash: fnv1a64:8b786fd96a7401e7
% canonical-lessons: RU-C347-odnako, RU-C347-zato, RU-C347-prichyom, RU-C347-itak, RU-C347-inymi-slovami

\chapter{However and moreover}
\label{ch:ru-however-and-moreover}
\hlchaptermodality{347}

\begin{chapteropening}
\textbf{By the end of this chapter:} \emph{I can link two ideas with however, but then, moreover, so, and in other words, and know what each one adds.}

Complete the last lesson of chapter 347: I can link two ideas with however, but then, moreover, so, and in other words, and know what each one adds.
\end{chapteropening}

\section[\texorpdfstring{\ru{однако} (odnáko)}{однако (odnáko)}]{\ru{однако} (odnáko) — however}

\label{lesson:RU-C347-odnako}



\begin{quote}
Before the new one: say the Russian for harm, then the Russian for to assess.
\end{quote}

\subsection*{You'll want to know: \ru{однако}}
\textbf{\ru{однако}} (\emph{odnáko}) — \textquotedblleft{}however\textquotedblright{}.

\begin{quote}
\ru{Было} \ru{холодно}, \ru{однако} \ru{мы} \ru{пошли} \ru{гулять}.
\end{quote}

\begin{quote}
\emph{Býlo khólodno, odnáko my poshlí gulyát.}
\end{quote}

\begin{quote}
It was cold; however, we went for a walk.
\end{quote}

Like \emph{\ru{но}}, but more formal. It can open the clause or sit inside it: \emph{\ru{Мы}, \ru{однако}, \ru{пошли} \ru{гулять}}.

\begin{cousinweb}
\textbf{\ru{однако}} is usually traced to \emph{\ru{один}} \textquotedblleft{}one\textquotedblright{}, with an old suffix; by that account it first meant \textquotedblleft{}in one and the same way, all the same\textquotedblright{}; \emph{\ru{один}} is inherited from Proto-Slavic.
\end{cousinweb}

\begin{grammarlens}[title={What you've built}]
One more word for telling, explaining and describing at B1.
\end{grammarlens}

\subsection*{Your turn}
\begin{itemize}
\raggedright
  \item \emph{Say it:} \emph{odnáko}
  \item \emph{Say it:} \emph{odnáko}, once more
  \item \emph{Say it:} a sentence of your own with \emph{odnáko}
  \item \emph{Recall:} say the Russian for harm, then the Russian for to assess, then say \emph{odnáko} again
\end{itemize}

\begin{tcolorbox}[breakable,colback=teal!4,colframe=teal!35!black,title={Before you move on}]
What does \ru{однако} mean? (\textquotedblleft{}However\textquotedblright{}.) Say it once more.
\end{tcolorbox}
\section[\texorpdfstring{\ru{зато} (zató)}{зато (zató)}]{\ru{зато} (zató) — but then, on the other hand}

\label{lesson:RU-C347-zato}



\begin{quote}
Before the new one: say the Russian for to assess, then the Russian for however.
\end{quote}

\subsection*{You'll want to know: \ru{зато}}
\textbf{\ru{зато}} (\emph{zató}) — \textquotedblleft{}but then, on the other hand\textquotedblright{}.

\begin{quote}
\ru{Квартира} \ru{маленькая}, \ru{зато} \ru{в} \ru{центре}.
\end{quote}

\begin{quote}
\emph{Kvartíra málenkaya, zató v tséntre.}
\end{quote}

\begin{quote}
The flat is small, but then it is in the centre.
\end{quote}

It brings in a plus that makes up for a minus: \emph{\ru{дорого}, \ru{зато} \ru{удобно}}. Often \emph{\ru{но} \ru{зато}}. Stress on the second syllable.

\begin{cousinweb}
\textbf{\ru{зато}} is fused from \emph{\ru{за}} \textquotedblleft{}for, in return for\textquotedblright{} and \emph{\ru{то}} \textquotedblleft{}that\textquotedblright{}: the good side you get in return for the bad.
\end{cousinweb}

\begin{grammarlens}[title={What you've built}]
One more word for telling, explaining and describing at B1.
\end{grammarlens}

\subsection*{Your turn}
\begin{itemize}
\raggedright
  \item \emph{Say it:} \emph{zató}
  \item \emph{Say it:} \emph{zató}, once more
  \item \emph{Say it:} a sentence of your own with \emph{zató}
  \item \emph{Recall:} say the Russian for to assess, then the Russian for however, then say \emph{zató} again
\end{itemize}

\begin{tcolorbox}[breakable,colback=teal!4,colframe=teal!35!black,title={Before you move on}]
What does \ru{зато} mean? (\textquotedblleft{}But then, on the other hand\textquotedblright{}.) Say it once more.
\end{tcolorbox}
\section[\texorpdfstring{\ru{причём} (prichyóm)}{причём (prichyóm)}]{\ru{причём} (prichyóm) — and what is more, moreover}

\label{lesson:RU-C347-prichyom}



\begin{quote}
Before the new one: say the Russian for however, then the Russian for but then.
\end{quote}

\subsection*{You'll want to know: \ru{причём}}
\textbf{\ru{причём}} (\emph{prichyóm}) — \textquotedblleft{}and what is more, moreover\textquotedblright{}.

\begin{quote}
\ru{Он} \ru{опоздал}, \ru{причём} \ru{уже} \ru{в} \ru{третий} \ru{раз}.
\end{quote}

\begin{quote}
\emph{On opozdál, prichyóm uzhé v trétiy raz.}
\end{quote}

\begin{quote}
He was late, and what is more, for the third time.
\end{quote}

It adds a detail that sharpens what came before. Stress on the second syllable.

\begin{cousinweb}
\textbf{\ru{причём}} is fused from \emph{\ru{при}} \textquotedblleft{}at, alongside\textquotedblright{} and \emph{\ru{чём}}, the prepositional case of \emph{\ru{что}}: \textquotedblleft{}alongside which\textquotedblright{}.
\end{cousinweb}

\begin{grammarlens}[title={What you've built}]
One more word for telling, explaining and describing at B1.
\end{grammarlens}

\subsection*{Your turn}
\begin{itemize}
\raggedright
  \item \emph{Say it:} \emph{prichyóm}
  \item \emph{Say it:} \emph{prichyóm}, once more
  \item \emph{Say it:} a sentence of your own with \emph{prichyóm}
  \item \emph{Recall:} say the Russian for however, then the Russian for but then, then say \emph{prichyóm} again
\end{itemize}

\begin{tcolorbox}[breakable,colback=teal!4,colframe=teal!35!black,title={Before you move on}]
What does \ru{причём} mean? (\textquotedblleft{}And what is more, moreover\textquotedblright{}.) Say it once more.
\end{tcolorbox}
\section[\texorpdfstring{\ru{итак} (iták)}{итак (iták)}]{\ru{итак} (iták) — so, and so (summing up)}

\label{lesson:RU-C347-itak}



\begin{quote}
Before the new one: say the Russian for but then, then the Russian for and what is more.
\end{quote}

\subsection*{You'll want to know: \ru{итак}}
\textbf{\ru{итак}} (\emph{iták}) — \textquotedblleft{}so, and so (summing up)\textquotedblright{}.

\begin{quote}
\ru{Итак}, \ru{все} \ru{готовы}? \ru{Тогда} \ru{начинаем}.
\end{quote}

\begin{quote}
\emph{Iták, vse gotóvy? Togdá nachináyem.}
\end{quote}

\begin{quote}
So, is everyone ready? Then we begin.
\end{quote}

It opens a summary or a conclusion. Stress on the second syllable.

\begin{cousinweb}
\textbf{\ru{итак}} is fused from \emph{\ru{и}} \textquotedblleft{}and\textquotedblright{} and \emph{\ru{так}} \textquotedblleft{}so\textquotedblright{}, both inherited Slavic words.
\end{cousinweb}

\begin{grammarlens}[title={What you've built}]
One more word for telling, explaining and describing at B1.
\end{grammarlens}

\subsection*{Your turn}
\begin{itemize}
\raggedright
  \item \emph{Say it:} \emph{iták}
  \item \emph{Say it:} \emph{iták}, once more
  \item \emph{Say it:} a sentence of your own with \emph{iták}
  \item \emph{Recall:} say the Russian for but then, then the Russian for and what is more, then say \emph{iták} again
\end{itemize}

\begin{tcolorbox}[breakable,colback=teal!4,colframe=teal!35!black,title={Before you move on}]
What does \ru{итак} mean? (\textquotedblleft{}So, and so (summing up)\textquotedblright{}.) Say it once more.
\end{tcolorbox}
\section[\texorpdfstring{\ru{иными} \ru{словами} (inými slovámi)}{иными словами (inými slovámi)}]{\ru{иными} \ru{словами} (inými slovámi) — in other words}

\label{lesson:RU-C347-inymi-slovami}



\begin{quote}
Before the new one: say the Russian for and what is more, then the Russian for so.
\end{quote}

\subsection*{You'll want to know: \ru{иными} \ru{словами}}
\textbf{\ru{иными} \ru{словами}} (\emph{inými slovámi}) — \textquotedblleft{}in other words\textquotedblright{}.

\begin{quote}
\ru{Он} \ru{не} \ru{ответил}, \ru{иными} \ru{словами}, \ru{он} \ru{против}.
\end{quote}

\begin{quote}
\emph{On ne otvétil, inými slovámi, on prótiv.}
\end{quote}

\begin{quote}
He did not answer; in other words, he is against it.
\end{quote}

The instrumental plural of \emph{\ru{иной}} \textquotedblleft{}other\textquotedblright{} and \emph{\ru{слово}} \textquotedblleft{}word\textquotedblright{}. \emph{\ru{Другими} \ru{словами}} means the same.

\begin{cousinweb}
\textbf{\ru{иными} \ru{словами}} is the instrumental plural of \emph{\ru{иной}} \textquotedblleft{}other\textquotedblright{}, the adjective behind \emph{\ru{иначе}}, and \emph{\ru{слово}} \textquotedblleft{}word\textquotedblright{}, inherited from Proto-Slavic \emph{*slovo} (EtymDB).
\end{cousinweb}

\begin{grammarlens}[title={What you've built}]
One more word for telling, explaining and describing at B1.
\end{grammarlens}

\subsection*{Your turn}
\begin{itemize}
\raggedright
  \item \emph{Say it:} \emph{inými slovámi}
  \item \emph{Say it:} \emph{inými slovámi}, once more
  \item \emph{Say it:} a sentence of your own with \emph{inými slovámi}
  \item \emph{Recall:} say the Russian for and what is more, then the Russian for so, then say \emph{inými slovámi} again
\end{itemize}

\begin{tcolorbox}[breakable,colback=teal!4,colframe=teal!35!black,title={Before you move on}]
What does \ru{иными} \ru{словами} mean? (\textquotedblleft{}In other words\textquotedblright{}.) Say it once more.
\end{tcolorbox}
